% %%%%%%%%%%%%%%%%%%%%%%%%%%%%%%%%%%%%%%%%%%%%%%%%%%%%%%%%%%%%%%%%%%%%%%%%%%%%%%%%%
%
%	cap5.tex 
%  -----------
%
%	Neste arquivo você escreve o capítulo. Deve sempre iniciar com o ambiente
%	\chapter{Nome do Capitulo}
%
% %%%%%%%%%%%%%%%%%%%%%%%%%%%%%%%%%%%%%%%%%%%%%%%%%%%%%%%%%%%%%%%%%%%%%%%%%%%%%%%%%
%


\chapter{Verificação do software}
% ----------------------------------------------------------

\section{Verificações parciais do modelo estrutural}

A estrutura utilizada no primeiro exemplo de verificação do modelo estrutural foi retirada do livro de \citeonline{Weaver2018} conforme pode se ver na Figura \ref{Fig.6.1}.

\begin{figure}[H]
    \centering
    \refstepcounter{figure}
    \label{Fig.6.1}

    \begin{minipage}{10cm}
        \fontsize{10pt}{12pt}\selectfont
        Figura \thefigure: Exemplo do livro \textit{Matrix Analysis of Framed Structures}
    \end{minipage}

    \includegraphics[scale=1.5]{exemplo1_livro.png}

    \begin{minipage}{10cm}
        \fontsize{10pt}{12pt}\selectfont
        Fonte: \citeonline{Weaver2018}
    \end{minipage}
\end{figure}

A estrutura foi modelada no Ftool, conforme apresentado na Figura \ref{Fig.6.2}. As variáveis necessárias para o cálculo foram definidas com os seguintes valores: módulo de elasticidade ($E = 10000,\text{kN/cm}^2$), área da seção transversal ($A = 10\text{cm}^2$), momento de inércia ($I = 1000,\text{cm}^4$), carga concentrada ($P = 10,\text{kN}$), comprimento ($L = 100,\text{cm}$), momento concentrado ($PL = 1000,\text{kN}\cdot\text{cm}$), carga distribuída ($2{,}4P/L = 0{,}24,\text{kN/cm}$), carga concentrada ($2P = 20,\text{kN}$) e distâncias ($3L/4 = 75,\text{cm}$ e $L/2 = 50,\text{cm}$).

\begin{figure}[H]
    \centering
    \refstepcounter{figure}
    \label{Fig.6.2}

    \begin{minipage}{10cm}
        \fontsize{10pt}{12pt}\selectfont
        Figura \thefigure: Exemplo gerado no Ftool
    \end{minipage}

    \includegraphics[scale=2]{exemplo1.png}

    \begin{minipage}{10cm}
        \fontsize{10pt}{12pt}\selectfont
        Fonte: Elaborado pelo autor
    \end{minipage}
\end{figure}

Os resultados obtidos pelo Ftool para valores de reações e diagramas de esforços solicitantes são apresentados na figura \ref{Fig.6.3} a seguir.

\begin{figure}[H]
    \centering
    \refstepcounter{figure}
    \label{Fig.6.3}

    \begin{minipage}{10cm}
        \fontsize{10pt}{12pt}\selectfont
        Figura \thefigure: Resultados obtidos pelo Ftool
    \end{minipage}

    \includegraphics[scale=0.3]{resultados_ftool.png}

    \begin{minipage}{10cm}
        \fontsize{10pt}{12pt}\selectfont
        Fonte: Elaborado pelo autor
    \end{minipage}
\end{figure}

\section{Verificações parciais da flexão simples}

Para verificar a rotina de dimensionamento à flexão simples, foi reproduzido o exemplo apresentado por \citeonline[p. 130--131]{carvalho2014calculo}, no qual se determina a área de armadura longitudinal necessária para uma seção retangular de concreto armado submetida a um momento fletor.

\begin{table}[H]
\centering
\caption{Dados do exemplo de referência}
\label{tab:exemplo_flexao}
\begin{tabular}{lc}
\hline
\textbf{Parâmetro} & \textbf{Valor} \\
\hline
Largura da seção ($b_w$) & 12 cm \\
Altura útil ($d$) & 29 cm \\
Resistência à compressão ($f_{ck}$) & 20 MPa \\
Resistência característica do aço ($f_{yk}$) & 500 MPa \\
Momento solicitante ($M$) & 12,2 kN$\cdot$m \\
\hline
\end{tabular}
\end{table}

Inicialmente, é preciso determinar os parâmetros de cálculo relativos ao concreto e ao aço, conforme as expressões:

\begin{align}
f_{cd} &= \frac{f_{ck}}{\gamma_c}
= \frac{20}{1,4}
= 14,286\,\text{MPa} \\
f_{yd} &= \frac{f_{yk}}{\gamma_s}
= \frac{500}{1,15}
= 434,78\,\text{MPa} \\
\varepsilon_{yd} &= \frac{f_{yd}}{E_s}
= 0,002070
\end{align}

Para o concreto, adotam-se $\alpha_c=0,85$, $\eta_c=1$ e $\lambda=0,8$. O limite entre os domínios 3 e 4 é inicialmente calculado por:

\begin{equation}
x_{lim}=
\frac{\varepsilon_{cu}}
{\varepsilon_{yd}+\varepsilon_{cu}}\,d
=0,1822\,\text{m}
\end{equation}

Como o valor obtido corresponde a $x_{lim}>0,45d$, adota-se $x_{lim}=0,45d=0,1305\,\text{m}$, conforme o critério estabelecido no fluxograma de dimensionamento.

A verificação da necessidade de armadura dupla é realizada a partir do momento limite da seção:

\begin{equation}
M_{dlim}=
\alpha_c\eta_c f_{cd}b_w\lambda x_{lim}
\left(d-\frac{\lambda x_{lim}}{2}\right)
\end{equation}

Substituindo os valores, obtém-se $M_{dlim}=36,18\,\text{kN}\cdot\text{m}$. Como $M_d\leq M_{dlim}$, a seção pode ser dimensionada com armadura simples, sendo considerada nula a armadura comprimida adicional.

A posição da linha neutra é determinada pela expressão:

\begin{equation}
x=
\frac{d-\sqrt{d^2-
\frac{2M_d}{\alpha_c\eta_c f_{cd}b_w}}}{\lambda}
\end{equation}

Resultando em $x=0,05464\,\text{m}$. A área de armadura longitudinal é então calculada por:

\begin{equation}
A_{s1}=
\frac{\alpha_c\eta_c f_{cd}b_w\lambda x}{f_{yd}}
=1,465\,\text{cm}^2
\end{equation}

O dimensionamento resulta em $A_{s1}=1,465\,\text{cm}^2$, valor próximo ao obtido no exemplo de referência, de $1,46\,\text{cm}^2$, evidenciando a concordância entre os resultados.

O segundo exemplo de verificação foi retirado de \citeonline[p.~132--133]{carvalho2014calculo} e corresponde ao mesmo problema apresentado no exemplo anterior, mantendo-se as dimensões da seção e o momento de cálculo, mas alterando-se a resistência característica à compressão do concreto para $f_{ck}=90\,\text{MPa}$. Assim, são mantidos $b_w=0,12\,\text{m}$, $d=0,29\,\text{m}$ e $M_d=17,08\,\text{kN}\cdot\text{m}$.

Inicialmente, são determinados os coeficientes relativos ao concreto, considerando os parâmetros correspondentes ao Grupo II:

\begin{align}
\lambda &= 0,8-\frac{f_{ck}-50}{400}=0,7 \\
\alpha_c &= 0,85\left[1-\frac{f_{ck}-50}{200}\right]=0,68
\end{align}

A resistência de cálculo do concreto é:

\begin{equation}
f_{cd}=\frac{90}{1,4}=64,286\,\text{MPa}
\end{equation}

A posição da linha neutra é obtida pela expressão:

\begin{equation}
x=\frac{d\pm\sqrt{d^2-
2\left(\frac{M_d}{b_w\alpha_cf_{cd}}\right)}}{\lambda}
\end{equation}

Substituindo os valores, obtêm-se duas soluções, $x_1=0,812\,\text{m}$ e $x_2=0,0164\,\text{m}$. Como a primeira solução posiciona a linha neutra fora da seção transversal, adota-se $x=0,0164\,\text{m}$.

Para verificar o domínio de deformação, determina-se a posição da linha neutra no limite entre os domínios 2 e 3:

\begin{equation}
x_{2-3}=0,259d=0,259\cdot0,29
=0,0751\,\text{m}
\end{equation}

Como $x<x_{2-3}$, a seção encontra-se no domínio 2, de modo que o aço atinge o escoamento e pode-se adotar $f_s=f_{yd}=43,478\,\text{kN/cm}^2$.

O braço de alavanca é calculado por:

\begin{equation}
z=d-\frac{\lambda x}{2}
=0,29-\frac{0,7\cdot0,0164}{2}
=0,284\,\text{m}
\end{equation}

Por fim, a área de armadura longitudinal necessária é:

\begin{equation}
A_s=\frac{M_d}{z f_{yd}}
=\frac{17,08}{0,284\cdot43,478}
=1,39\,\text{cm}^2
\end{equation}




